\section{Local Slack and Macroscopic Averages}\label{f:section}
The estimates in this section provide the spatial compactness and moment
bounds used in the comparison argument. The principal probabilistic input
is the controlled-gradient theorem of Cohen--Alloro and Peled \cite{CAP}.
The local convex-volume argument uses integral log-concavity
\cite[Theorems 3 and 6]{Prekopa}. Its quantitative facet and block
estimates are proved below for the present triangular lattice.

\subsection{The finite-volume laws and the moment statement}
We identify the triangular lattice with $\mathbb Z^2$, with edges in the
directions $\pm e_1,\pm e_2,\pm(e_1-e_2)$. Distances are graph distances.
Let $V$ be a finite set of free vertices, and prescribe height zero on its
exterior vertex boundary and on any additional pins. Write $\mu_0$ for
uniform Lebesgue measure on the remaining heights satisfying the edge
constraints. A zero-pinned sample is extended by zero to the whole lattice.
Its defining polytope is compact and has positive volume on its free
coordinates. A bounded-data law $\mu_b$ has the same fixed set and prescribed
values $b$ with $|b|\le B$, and is considered only when its polytope has
positive volume. Constants below are independent of the number and shape
of the components of $V$.

\begin{proposition}[Macroscopic exponential moments]\label{f:averages}
Fix $c,C>0$ and $B<\infty$. Let $R\ge1$ and let deterministic real weights
$q_x$ obey
\begin{equation}\label{f:weights}
 \sum_x|q_x|\le C,\qquad \max_x|q_x|\le CR^{-2},\qquad
 \operatorname{supp}q\subset Q,
 \quad \operatorname{diam}Q\le CR,
\end{equation}
where $Q$ is an ambient coordinate box. Assume either $\sum_xq_x=0$, or
there is a connected set of pins $K$, of diameter at least $cR$, such that
$K\cup Q$ lies in a coordinate box of diameter at most $CR$. For a sample
$h$ from $\mu_b$, there are $a>0$ and $M<\infty$, depending only on
$c,C,B$, such that
\begin{equation}\label{f:average-mgf}
 \E\exp\left(a\left|\sum_xq_xh(x)\right|\right)\le M.
\end{equation}
The conclusion is uniform over mixtures of bounded pin values whenever
conditioning on those values leaves exactly the specified uniform law.
\end{proposition}

The zero-mass alternative covers differences of comparable-scale averages
in a free ball even when the nearest pin is much farther away. An ordinary
ball average, or a smooth average of that scale, satisfies
\eqref{f:weights}. Restriction of its original positive weights to the
domain also satisfies that condition. Renormalization after restriction
requires a new verification of the bound $\max|q_x|\le CR^{-2}$.
We first prove the proposition for zero pins. The proof will also identify
the exact deterministic estimates required of later comparison fields.

\subsection{A local filling-volume deficit}
Fix a large even integer $L$ and tile the lattice by
$Q_z=Lz+\{0,\ldots,L\}^2$. The interiors of distinct blocks have no
connecting edge. For shell data $s$ on $\partial Q$, let $K_Q(s)$ be the
polytope of fillings and put
\[
 F_Q(s)=\log\operatorname{vol}_{(L-1)^2}K_Q(s),
\]
with value $-\infty$ when the volume is zero. The inclusion
$\frac12(K_Q(s_1)+K_Q(s_2))\subset K_Q((s_1+s_2)/2)$ and
Brunn--Minkowski imply concavity of $F_Q$.

Concavity gives a nonnegative midpoint loss. We next show that, on blocks
with slack, this loss controls differences between the perturbations of
the shell heights at the four side midpoints. It cannot control a common
height shift, which leaves filling volume unchanged. The later coarse-grid
argument combines these differences; the connected pins or the zero-mass
hypothesis in Proposition~\ref{f:averages} remove the additive constant.

For each of the four sides, let $v$ be its midpoint and let $r$ be
the coordinate unit step pointing toward the block center. The inward
path is $v+kr$, $0\le k\le L/2$. Of the two tangential unit steps
on that side, let $t$ be the one with $r\cdot t=1/2$. The
\emph{compressed ray} is $v+\ell t$, $0\le\ell\le L/2$:
its graph distance to $v+kr$ is $\max(k,\ell)$. These choices define
the four tests below. Figure~\ref{fig:triangular} illustrates one of them.

For fixed $0<\delta\le\delta_0\le1/16$, an integer $k_0\ge4/\delta_0$,
and $s_0>0$, call a block \emph{convex-good} for $h$ if
\begin{align}
 |h(x)-h(y)|&\le1-\delta &&(x\sim y\text{ in }Q),\label{f:near}\\
 |h(x)-h(y)|&\le(1-\delta_0)\operatorname{dist}(x,y)
 &&(x,y\in Q,\ \operatorname{dist}(x,y)\ge k_0),\label{f:far}
\end{align}
and the following four non-tie conditions hold. For each face center $v$,
take the radial coordinate path from $v$ to the block center. Exactly one
ray on its starting face has compressed triangular distances to that path;
if $b_1$ is its first vertex, require $|h(b_1)-h(v)|\ge s_0$.
Take $L\ge8(k_0+2)/\delta_0$.

\begin{figure}[ht]
\centering
\begin{tikzpicture}[scale=1.05,>=Latex]
\foreach \i in {0,...,4} {
  \foreach \j in {0,...,3} {
    \coordinate (v\i\j) at ({\i+0.5*\j},{0.8660254*\j});
    \fill[gray] (v\i\j) circle (1.1pt);
  }
}
\foreach \i in {0,...,3} {
  \foreach \j in {0,...,3} {
    \draw[gray!35] ({\i+0.5*\j},{0.8660254*\j}) --
      ({\i+1+0.5*\j},{0.8660254*\j});
  }
}
\foreach \i in {0,...,4} {
  \foreach \j in {0,...,2} {
    \draw[gray!35] ({\i+0.5*\j},{0.8660254*\j}) --
      ({\i+0.5*(\j+1)},{0.8660254*(\j+1)});
  }
}
\foreach \i in {0,...,3} {
  \foreach \j in {0,...,2} {
    \draw[gray!35] ({\i+1+0.5*\j},{0.8660254*\j}) --
      ({\i+0.5*(\j+1)},{0.8660254*(\j+1)});
  }
}
\draw[blue!65!black,very thick] (v00)--(v03);
\draw[red!65!black,very thick] (v20)--(v30);
\draw[green!45!black,thick,dashed] (v20)--(v11)--(v02)--(v03);
\draw[orange!80!black,thick,dotted] (v30)--(v21)--(v12)--(v03);
\node[below left] at (v00) {$v$};
\node[left] at (v01) {$v+e_2$};
\node[above] at (v03) {$v+3e_2$};
\node[below] at (v20) {$x_2$};
\node[below] at (v30) {$x_3$};
\draw[->] (4.5,0.35)--(5.2,0.35) node[right] {$e_1$};
\draw[->] (4.5,0.35)--(4.85,0.956) node[above right] {$e_2$};
\end{tikzpicture}
\caption{The compressed ray in the triangular exposure argument. Here
$x_k=v+ke_1$, and the highlighted test edge is $x_2x_3$. Both drawn paths
from its endpoints to $v+3e_2$ have length three. For a ray vertex
$v+\ell e_2$, the two distances are $\max(k,\ell)$ and
$\max(k+1,\ell)$, rather than
$k+\ell$ and $k+1+\ell$. The sign of
$h(v+e_2)-h(v)$ selects which facet can be exposed with uniform slack;
its near-zero case must be excluded in that argument.}
\label{fig:triangular}
\end{figure}



\begin{lemma}[Triangular block deficit]\label{f:block}
For these fixed parameters there is $c_*>0$ such that, if $Q$ is
convex-good for $h$, $s=h|_{\partial Q}$,
$\chi:\partial Q\to\R$, and both $K_Q(s+\chi)$ and
$K_Q(s-\chi)$ have positive volume, then every two face centers $v,w$
satisfy
\begin{equation}\label{f:block-deficit}
 g_Q:=F_Q(s)-\tfrac12\bigl(F_Q(s+\chi)+F_Q(s-\chi)\bigr)
 \ge c_*|\chi(v)-\chi(w)|^2.
\end{equation}
\end{lemma}
\begin{proof}
We first expose one sign of each radial edge with uniform slack in all
other constraints. Write $x=v+ke_1$, $y=v+(k+1)e_1$ after applying a lattice
symmetry. If $k\ge k_0$, prescribe $a(x)=h(x)$ and $a(y)=h(x)+1$.
Both points are at distance at least $k_0$ from the shell, and the change
at $y$ is at most two. Hence, for sufficiently small $\beta\le\delta/2$,
\[
 |a(z)-h(b)|\le(1-\delta_0)\operatorname{dist}(z,b)+2
 \le(1-\beta)\operatorname{dist}(z,b)
 \quad(z=x,y,\ b\in\partial Q).
\]
If $k<k_0$, put $b_l=v+le_2$ on the compressed ray and
$f_l=h(b_l)-h(v)$. Its distances are
\[
 \operatorname{dist}(x,b_l)=\max(k,l),\qquad
 \operatorname{dist}(y,b_l)=\max(k+1,l).
\]
Choose $\tau=\operatorname{sign}f_1$, and prescribe
$a(x)=h(v)$ and $a(y)=h(v)+\tau$. For $l\ge1$,
\[
 s_0-(1-\delta)(l-1)\le\tau f_l\le(1-\delta)l,
\]
so
\begin{equation}\label{f:triangular-slack}
 |\tau-f_l|\le l-\min\{s_0+\delta(l-1),1+\delta l\}
 \le(1-\min\{s_0,\delta\})l.
\end{equation}
The difference at $x$ is bounded by $(1-\delta)l$. On the opposite ray,
the distances are $k+l,k+1+l$. The estimate
$1+(1-\delta)l\le(1-\delta/2)(k+1+l)$ holds when $k+l\ge1$.
The omitted case is the designated edge itself. For another face, put
$l=\operatorname{dist}(v,b)\ge L/2$. Both prescribed differences are at
most $1+(1-\delta_0)l$ and both distances are at least $l-k-1$.
The choice of $L$ gives the required relative slack there as well.

Choose $\beta<\min\{\delta/2,s_0/2,1/4\}$, decreasing it if necessary.
Give the designated edge length one and all other edges length $1-\beta$.
All prescribed pairs are Lipschitz for this weighted metric. Its distance
across the designated edge is one, since every alternative path has at
least two edges and $2(1-\beta)>1$. Extending by the minimum of the upper
distance cones gives a filling with designated gradient $\tau$ and all
other gradients at most $1-\beta$ in absolute value.

We need the following quantitative finite-dimensional consequence of
slice log-concavity, and give its proof.
If positive-volume convex bodies $K_-,K_+\subset\R^m$
have projection widths at most $A$ in a unit direction $u$, their top
facet areas divided by volume are at least $c_0>0$, and
$K\supset(K_-+K_+)/2$, then
\begin{equation}\label{f:facet-input}
 \log\frac{|K|}{\sqrt{|K_-||K_+|}}
 \ge\kappa\bigl[(b_+-u\cdot m_+)-(b_--u\cdot m_-)\bigr]^2.
\end{equation}
Here $m_\pm$ are the centroids, $b_\pm$ the top projection endpoints,
and $\kappa>0$ depends only on $A,c_0$. This is a finite-dimensional lemma;
it has no restriction on the spatial dimension of the lattice.

To prove \eqref{f:facet-input}, let $f_i$ be the normalized density
of the projection of uniform measure on $K_i$, and let $q_i(p)$
be its quantile, $0<p<1$, for $i\in\{-,+\}$. Parallel-slice
Brunn--Minkowski makes $f_i$ log-concave, with $f_i(b_i)\ge c_0$.
If $f_i(q_i(p))<c_0$, every density value to its left is at most
$f_i(q_i(p))$: otherwise log-concavity between that value and the
top slice gives a contradiction. Integration over an interval of
length at most $A$ then gives $p\le A f_i(q_i(p))$.
In the other case $q_i'(p)=1/f_i(q_i(p))\le c_0^{-1}$.
Thus, with $C=\max\{A,c_0^{-1}\}$,
\[
 0<q_i'(p)\le C/p.
\]
The midpoint slice of $K$ at $(q_-(p)+q_+(p))/2$ contains the
midpoint sum of the two corresponding slices. Slice
Brunn--Minkowski and change of variable along these increasing
levels yield
\[
 \frac{|K|}{\sqrt{|K_-||K_+|}}
 \ge\int_0^1\frac{q_-'+q_+'}{2\sqrt{q_-'q_+'}}\dd p
 \ge1+\frac1{8C^2}\int_0^1p^2(q_-'-q_+')^2\dd p.
\]
For the second inequality, subtract one from the integrand; its
numerator is $(q_-'-q_+')^2$ and its denominator is
$2\sqrt{q_-'q_+'}(\sqrt{q_-'}+\sqrt{q_+'})^2\le8C^2/p^2$.
Tonelli gives
$b_i-u\cdot m_i=\int_0^1p q_i'(p)\dd p$.
Cauchy--Schwarz therefore bounds the last integral below by the
squared difference $d^2$ of the centroid gaps. Since $|d|\le A$,
$\log(1+x)\ge x/(1+x)$ proves \eqref{f:facet-input} with
$\kappa=(8C^2+A^2)^{-1}$. All slice integrals may first be taken
over compact subintervals of $(0,1)$ and then by monotone limits.

Apply it to $K_\pm=K_Q(s\pm\alpha\chi)$ with
$\alpha\le\beta/(8L)$. Endpoint feasibility implies
$|\chi(b)-\chi(c)|\le\operatorname{dist}(b,c)\le2L$ on the shell.
Shift every interior coordinate in the exposed filling by
$\pm\alpha\chi(v)$ and use the new shell data. Every undesignated edge
retains slack at least $\beta/2$. Thus each selected facet contains a
Euclidean ball of radius $\beta/8$ in its hyperplane. Body volumes are at
most $(4L)^{(L-1)^2}$, and the projection widths are at most two.
The hypotheses of \eqref{f:facet-input} therefore hold uniformly.

Define $U=m_+-m_-$ on the interior and $U=2\alpha\chi$ on the shell.
For an interior radial edge use the direction
$\tau(e_y-e_x)/\sqrt2$. The top projection endpoints are both $1/\sqrt2$,
so the difference of the two centroid gaps in~\eqref{f:facet-input} is
$-\tau(U(y)-U(x))/\sqrt2$. For an edge from a face center $v$ to an
interior vertex $y$, use the direction $\tau e_y$. Its top endpoints are
$1+\tau(s(v)\pm\alpha\chi(v))$, so the difference of the centroid gaps
is $-\tau(U(y)-U(v))$. Thus~\eqref{f:facet-input} bounds every
$|U(y)-U(x)|^2$ by a fixed multiple of the common midpoint gap $g_\alpha$.
The bodies and their centroids are the same for all radial edges; varying
$\tau$ changes no squared increment. Telescoping at most $L$ radial edges
between $v,w$ and using $g_\alpha\le g_Q$ proves the lemma.
\end{proof}

When $f_1=0$ and $k=0$, both choices $a(y)=h(v)\pm1$ also saturate the
edge $yb_1$. The non-tie requirement above repairs this actual triangular
obstruction; it is not supplied by the square-lattice facet argument.

\subsection{Sparse exceptions and short stationary steps}
\begin{lemma}[Zero-pinned sparse blocks]\label{f:sparse}
For every $\rho>0$ the fixed parameters above and a fixed enlargement
$Q_z^*$ can be chosen so that the following holds uniformly over zero-pinned
finite volumes. Declare every block within a fixed block distance of a pin
to be a pin block. Among the remaining blocks declare a block bad when it
fails \eqref{f:near}, the long-distance condition \eqref{f:far} on $Q_z^*$,
or one of the four non-tie tests. Then
\begin{equation}\label{f:animal}
 \Prob(A\text{ is entirely bad})\le\rho^{|A|}
\end{equation}
for every finite set $A$ of block indices. The collar defining pin blocks
contains all enlargements $Q_z^*$. Fixed finite thickenings of the bad set
also satisfy such a bound with an arbitrarily small parameter, after
adjusting the initial parameters.
\end{lemma}
\begin{proof}
The controlled-gradient theorem \cite[Theorem 1.1]{CAP}, on a graph of
maximum degree six with an arbitrary nonempty zero fixed set, states that
for distinct edges $e_1,\ldots,e_j$ and $0<u\le1/8$,
\[
 \Prob\bigl(|\nabla_{e_i}h|\ge1-u\text{ for all }i\bigr)
 \le(C_6u)^{j/C_6}.
\]
Fix one shortest path for every pair in an enlarged block. A path of
length $m$ violating linear slack with $\delta_0=\eps/2$ has at least
$m/2$ gradients of magnitude at least $1-\eps$. Choose $\eps$ sufficiently
small to dominate the $2^m$ subsets of this path, then $k_0$ large enough
to dominate the polynomial number of pairs. Choose $L$ at a fixed
sufficiently large multiple of $k_0/\delta_0$, and finally choose $\delta$
small enough to dominate the finite edge count. For disjoint enlargements,
the edge witnesses are distinct. Apply the controlled-gradient bound to
the short-edge and long-path witnesses separately and bound the intersection
by the geometric mean of these two bounds. This gives an arbitrarily small
product bound for failures of the first two tests.

For a non-tie test, condition on its enlarged shell. If there is a filling
with the stated long slack, fix one such filling and write $x,y$ for the
tested edge. Prescribe its values to be $h(x),h(x)+1/2$, or
$h(x),h(x)-1/2$. The distance of the edge to the shell is a sufficiently
large fixed multiple of $1/\delta_0$; long slack and the triangle inequality
make both prescriptions compatible with the shell. Upper-cone extension
produces two feasible fillings. Contracting the entire conditional body
toward either filling with ratio $1/4$ gives a volume fraction $4^{-m}$
in each of the difference intervals $[1/8,5/8]$ and $[-5/8,-1/8]$,
where $m$ is the fixed number of free coordinates. The difference marginal
is log-concave. A one-dimensional log-concave density placing a fixed
positive mass in these separated intervals has bounded supremum: its
half-height points lie within $2/\sup f$ of a mode, and concavity of
$\log f$ bounds the farther tails exponentially. Consequently
\[
 \Prob\bigl(|h(x)-h(y)|<s_0,\ \text{long slack}\mid\text{shell}\bigr)
 \le C_Ls_0.
\]
If no long-slack filling exists the event is empty. Conditional laws in
disjoint enlarged interiors are independent given their complement.
Choosing $s_0$ small, combining the possible witnesses, and coloring the
block grid by finitely many residue classes proves \eqref{f:animal}.
For a fixed thickening, extract a separated positive fraction of witnesses
in the original set; the same counting proves the last assertion.
\end{proof}

\begin{lemma}[Stationary step displacement]\label{f:step}
Let $(H_t)_{0\le t\le T}$ be a stationary reversible normally reflected
diffusion in a finite convex polytope, with Brownian covariance $2I$ and
invariant uniform law. The same assertion holds with constant drift $b$
and invariant density proportional to $e^{b\cdot h}$.
For distinct free coordinates $x_1,\ldots,x_j$,
\begin{equation}\label{f:step-tail}
 \Prob\left(\sup_{t\le T}|H_t(x_i)-H_0(x_i)|>r
                  \text{ for all }i\right)
 \le\bigl(Ce^{-r^2/(CT)}\bigr)^{j/2}.
\end{equation}
Thus the zero-pinned good-block conditions persist with relaxed fixed
constants throughout a sufficiently short step, apart from arbitrarily
sparse exceptional blocks. The same implication holds for any other
stationary constrained law for which the initial block tests have the
joint sparse bound.
\end{lemma}
\begin{proof}
Normal reflection has a forward Brownian martingale part and a
finite-variation part consisting of boundary reflection and any constant
drift. Reversibility gives the same decomposition backwards
from $T$. Adding the forward increment and the negative reversed increment
cancels the finite-variation terms, so each coordinate displacement is half the
sum of the corresponding forward and reversed martingale increments.
Each of the two martingale vectors has independent variance-two Brownian
coordinates, although the two vectors need not be independent. If all $j$
coordinate displacements exceed $r$, at least $j/2$ are witnessed in one
of the two directions. A union over subsets and the Brownian maximal
inequality prove \eqref{f:step-tail}. A failed block persistence test is
witnessed by a coordinate displacement exceeding a fixed constant.
Separated blocks have distinct witnesses; their finite coordinate count
is absorbed by taking $T$ smaller. Edge slack, long-distance slack, and
non-ties all persist once their two endpoint displacements are below
the appropriate fixed fraction of their original margins.
\end{proof}

\subsection{Filling holes and testing macroscopic averages}
We next isolate the deterministic geometry. It is valid for comparison
fields as well as the displacements of shell data used below.
\begin{lemma}[Sparse-hole extension]\label{f:hole}
Let a random bad set on the square grid satisfy \eqref{f:animal} with
sufficiently small parameter. Fill the bounded complementary components
of each bad star-connected cluster and retain its outermost filled hulls.
A fixed layer of enlargement is allowed. There is an extension procedure
which, for a field $u$ on good sites, fills each hole by its value at one
surrounding good site. Here good sites lie outside the filled holes and
good edges join neighboring good sites. Write $v$ for the resulting field.
A hole is incident to $z$ if it contains $z$ or a neighbor of $z$;
let $d_z$ be the largest diameter of an incident hole, zero if none,
and put $d_e=\max(d_x,d_y)$ for $e=\{x,y\}$. Uniformly in deterministic $z$,
\begin{equation}\label{f:hole-tail}
 \Prob(d_z\ge r)\le Ce^{-cr}.
\end{equation}
For every $1<p<2$, on a square of side $N$ with a fixed larger buffer,
\begin{equation}\label{f:hole-weighted}
 \|\nabla v\|_{\ell^p}
 \le E^{1/2}
 \left(\sum_e W_e^{2/(2-p)}a_e^{-p/(2-p)}\right)^{(2-p)/(2p)},
 \quad W_e\le C_p(1+d_e)^{2p},
\end{equation}
where $E=\sum_{e\text{ good}}a_e|\nabla u(e)|^2$, $a_e>0$.
In particular, bounded spatial moments of sufficiently high powers of
$d_e$ and $a_e^{-1}$ imply
\begin{equation}\label{f:hole-scaled}
 \|\nabla v\|_{\ell^p}\le C N^{2/p-1}\sqrt E.
\end{equation}
If $u$ is the restriction of a globally $K$-Lipschitz field on the grid,
then $|u(z)-v(z)|\le CK(1+d_z)$.

For every fixed $q,D<\infty$, the spatial average of $(1+d_z)^q$ in a
square of side $N$ is bounded by a deterministic constant outside an event
of probability $O(N^{-D})$. In a box of side $n$, these bounds hold
simultaneously on all subboxes of radius at least $n^{1/2}$, after altering
the failure exponent. Filling can be chosen equivariantly under translations
of the block grid. No reciprocal-weight assertion follows from sparsity
alone; those moments are separate hypotheses in \eqref{f:hole-weighted}.
\end{lemma}
\begin{proof}
Follow the filled outer grid contour of a hole and its exterior good rim.
After a fixed thickening this gives nearest-neighbor rim paths. A hole
of diameter $d$ has at most $C(d+1)^2$ vertices and boundary edges, and
any required rim path has at most that length. Choose the lexicographically
first rim site as the filling value. For an edge entering a hole, telescope
along its rim; the inequality $|\sum_{i=1}^m t_i|^p\le
m^{p-1}\sum_i|t_i|^p$, followed by the boundary-edge count, gives
\[
 \sum|\nabla v|^p\le\sum_{e\text{ good}}W_e|\nabla u(e)|^p,
 \qquad W_e\le C_p(1+d_e)^{2p}.
\]
Outer rims have bounded overlap. Holder with exponents $2/p$ and
$2/(2-p)$ gives \eqref{f:hole-weighted}. For unweighted energy the required
hole exponent is $4p/(2-p)$; in the weighted case another Holder inequality
separates the hole and reciprocal-weight moments. A hole filling value lies
within distance $C(d+1)$ of each of its sites, proving the pointwise error.

If a deterministic site belongs to a hull of diameter at least $r$, a bad
star-animal of size $m\ge cr$ has a vertex within $Cm$ of it. There are
$O(m^2)$ root choices and at most $C_0^m$ such animals. Applying
\eqref{f:animal}, or extracting a fixed-fraction separated subset before
doing so, proves \eqref{f:hole-tail}. In particular all fixed moments
have uniformly bounded expectations.

Here is the stronger spatial statement. Clusters with size in $[k,2k)$
contribute at most $Ck^{q+2}$ each to the sum of hole moments. For $r$
disjoint such cluster witnesses in an $O(N)$ box, root and animal counting
and \eqref{f:animal} bound the probability by
\[
 \frac{(CN^2)^r}{r!}C_0^{2kr}\rho^{ckr}
 \le\left(C\frac{N^2}{r}e^{-c'k}\right)^r.
\]
For all sufficiently large fixed $k$, take
$r=\lceil N^2k^{-(q+4)}\rceil$. For $k\le A\log N$ and large $N$, the
resulting bound is
\[
 \exp(-cN^2/k^{q+3}).
\]
Larger clusters occur with
probability at most $CN^2e^{-cA\log N}$, which is $O(N^{-D})$ by
choosing $A$ large. On all dyadic size ranges the resulting contribution
to $N^{-2}\sum_z(1+d_z)^q$ is summable; clusters below the fixed cutoff
give a fixed deterministic bound. This proves the assertion. The same
argument excludes holes crossing the macroscopic buffer, and a union over
the polynomially many subboxes proves the stated simultaneous version.
Lexicographic selection commutes with translation; a random grid phase
can be adjoined when invariance under single lattice steps is desired.
\end{proof}

\begin{lemma}[Discrete trace against a spread measure]\label{f:trace-average}
Let $\nu$ be a positive measure of bounded mass in a square of side $O(N)$,
and assume $\nu(B(x,r))\le C((r+1)/N)^\alpha$ for $0\le r\le CN$.
If $1<p<2$ and $\alpha>2-p$, then
\[
 |\nu(v)-\nu(1)\overline v|
 \le C N^{1-2/p}\|\nabla v\|_{\ell^p},
\]
where $\overline v$ is the ordinary average over that square.
\end{lemma}
\begin{proof}
Discrete $L^1$ Poincare on successive dyadic squares, clipped to the
ambient square at its boundary, gives
\[
 |v(x)-\overline v|\le C\sum_e
       \frac{|\nabla v(e)|}{1+\operatorname{dist}(x,e)}.
\]
For $p'=p/(p-1)$, a dyadic contribution to the averaged kernel is
$F_r(e)=Cr^{-1}\nu(B(e,Cr))$. Its supremum is at most
$CN^{-\alpha}r^{\alpha-1}$ and its $\ell^1$ norm at most $Cr$.
Thus
\[
 \|F_r\|_{\ell^{p'}}\le
 C N^{-\alpha/p}r^{1+(\alpha-2)/p}.
\]
The dyadic series is bounded by $CN^{1-2/p}$ precisely under the stated
strict inequality on $\alpha$. Holder finishes the proof.
\end{proof}

A connected set of coarse pins of diameter at least $cN$ carries a
deterministic probability measure $\nu$ with the preceding ball bound for
$\alpha=1$: choose a coordinate projection of length at least $c'N$ and
one pin site over each attained integer coordinate. Connectivity ensures
there are no missing intermediate coordinates. Equal weights on the chosen
sites have ball masses at most $C(r+1)/N$.

\begin{lemma}[Uniform average coercivity]\label{f:average-coercivity}
Suppose $X$ is globally $K$-Lipschitz on the block grid, with fixed $K$.
Suppose either $u=X$, or $u=T_M(X)$ for fixed
$T_M(t)=\operatorname{sign}(t)(|t|-M)_+$. Let
$E=\sum_{e\text{ good}}|\nabla u(e)|^2$. For deterministic weights $w$
supported in an $O(N)$ square with $\sum|w|\le C$ and
$\max|w|\le CN^{-2}$, assume either $\sum w=0$, or $u=0$ on a connected
deterministic coarse pin set of diameter comparable to $N$ at distance
$O(N)$. There is an event of probability at least $1/2$, determined only
by the bad set, on which simultaneously for every such field $X$,
\begin{equation}\label{f:coercivity}
 \left|\sum_z w_zX_z\right|\le C_0+C_1\sqrt E.
\end{equation}
\end{lemma}
\begin{proof}
Use Lemma~\ref{f:hole} with $p=3/2$. A Markov bound on the twelfth spatial
hole moment, and exclusion of holes crossing a fixed buffer, give
$\|\nabla v\|_{3/2}\le CN^{1/3}\sqrt E$. The filling error tested against
$|w|$ is bounded by the same spatial moment. In the anchored case choose
the deterministic pin measure $\nu$ just constructed. Although filling
may overwrite a pin value, the Lipschitz bound gives
\[
 |\nu(v)|\le CK\sum_z\nu(z)(1+d_z).
\]
The right side has bounded expectation by \eqref{f:hole-tail}; another
Markov bound makes it a fixed constant. Applying
Lemma~\ref{f:trace-average} to $\nu$ and $|w|$ proves
\eqref{f:coercivity}. In the zero-mass case the ordinary mean cancels and
no pin measure is needed. The truncation error is at most $M\sum|w|$.
Choose the two moment thresholds so their failure probabilities sum to
at most $1/4$; the buffer event has failure at most $1/4$ for large $N$.
Bounded $N$ is absorbed by deterministic Lipschitzness and the anchor or
zero-mass condition. All error bounds depend on hole sizes rather than
on $X$, which proves the simultaneous assertion.
\end{proof}

This last lemma is an unscaled estimate. Its bounded additive error on
the pin set does not assert convergence to zero Sobolev trace in a later
rescaled limit.

\subsection{From deficits to an exponential moment}
\begin{proof}[Proof of Proposition~\ref{f:averages} for zero pins]
For $R/L$ large integrate the interiors of pin-free blocks in a fixed
large box around the statistic and its anchor. Retain all other vertices,
including every vertex of pin blocks, and write $s$ for their heights.
The unnormalized marginal density is exactly
\begin{equation}\label{f:product}
 p(s)=\mathbf1_{\mathcal A}(s)\prod_Q e^{F_Q(s|_{\partial Q})},
\end{equation}
where $\mathcal A$ imposes the remaining edge constraints. This density is
even, log-concave, compactly supported, and upper semicontinuous. Let
$Z=\int p>0$.

For feasible $s\pm\chi$ with positive density, put $X_z=\chi(v_z)$ at
a fixed face-center representative of each block. The difference of any
two full endpoint fillings, divided by two, extends $\chi$ as a
1-Lipschitz field vanishing on pins. Consequently neighboring $X$ values
differ by at most $CL$. Choose $M$ above the distance of every representative
and shared face center in or next to a pin block to a pin, and use
$u=T_M(X)$. It vanishes on pin blocks. For neighboring ordinary good blocks
$z,w$, their common face center and Lemma~\ref{f:block} give
\[
 |u_z-u_w|^2\le2c_*^{-1}(g_z+g_w).
\]
For a good block next to a pin block, the common face center $f$ satisfies
$|\chi(f)|\le M$, so $|u_z|^2\le|\chi(v_z)-\chi(f)|^2\le g_z/c_*$.
Treating pin blocks as good zero blocks therefore gives
\begin{equation}\label{f:energy-deficits}
 \sum_{e\text{ good}}|\nabla u(e)|^2\le C\sum_Q g_Q.
\end{equation}

Partition vertices into half-open blocks and replace the statistic by
the retained linear functional $\ell(s)=\sum_z w_zs(v_z)$, where
$w_z=\sum_{x\text{ in block }z}q_x$. On every original configuration,
\begin{equation}\label{f:representative-error}
 \left|\sum_xq_xh(x)-\ell(s)\right|\le CL\sum_x|q_x|.
\end{equation}
The coarse weights satisfy the hypotheses of
Lemma~\ref{f:average-coercivity}, and exact zero mass is preserved.
That lemma and \eqref{f:energy-deficits} show, on a full-configuration
event of probability at least $1/2$, simultaneously for every feasible
$\chi$,
\[
 |\ell(\chi)|\le C_0+C_2\sqrt{\sum_Qg_Q}.
\]
Take fixed $t>C_0$ and $D=((t-C_0)/C_2)^2>0$. Since the logarithmic
midpoint deficit of \eqref{f:product} is $\sum_Qg_Q$, it follows that
\begin{equation}\label{f:pointwise-overlap}
 \sqrt{p(s+\chi)p(s-\chi)}\le e^{-D}p(s)
 \quad\text{whenever }\ell(\chi)=t.
\end{equation}

Define $A_t(s)=\sup_{\ell(\chi)=t}\sqrt{p(s+\chi)p(s-\chi)}$.
Compactness of the support restricts the supremum to a compact set;
upper semicontinuity gives measurability. Log-concavity implies
$0\le A_t\le p$. The good event may depend on discarded midpoint
coordinates: the simultaneous assertion \eqref{f:pointwise-overlap}
still holds for its retained vector. Thus
\[
 Z^{-1}\int A_t\le1-\rho_0,\qquad
 \rho_0=(1-e^{-D})/2>0.
\]
If $\ell$ is zero on the free coordinates there is nothing to prove.
Otherwise choose linear coordinates $s=(r,\sigma)$ with $r=\ell(s)$,
absorbing their constant Jacobian into the complementary measure.
For the density $a(r)$ of $\ell(S)$, Pr\'ekopa--Leindler in $\sigma$
\cite[Theorems 3 and 6]{Prekopa} gives
\[
 \sqrt{a(r-t)a(r+t)}\le Z^{-1}\int A_t(r,\sigma)\dd\sigma.
\]
Integrating yields overlap at most $1-\rho_0$. Since $a$ is even and
log-concave, this overlap is at least
$\int\min\{a(r-t),a(r+t)\}\dd r=1-\int_{-t}^t a$.
Hence $\Prob(|\ell(S)|\le t)\ge\rho_0$. Its survival function is
log-concave, again by Prekopa--Leindler, so
\[
 \Prob(|\ell(S)|>jt)\le(1-\rho_0)^j\quad(j\ge1).
\]
This gives a positive uniform exponential moment. Equation
\eqref{f:representative-error} transfers it to the original statistic.
For bounded $R/L$, distance to the anchor, or cancellation of $q$, gives
the required deterministic bound.
\end{proof}

\subsection{Bounded values, sign constraints, and small tilts}
\begin{lemma}[One-sample bounded-data comparison]\label{f:bounded}
There is a coupling of $h_b\sim\mu_b$ and one sample $h_0\sim\mu_0$ with
$|h_b-h_0|\le B$ at every vertex, whenever $|b|\le B$.
Consequently Proposition~\ref{f:averages} holds for bounded data and their
admissible mixtures.
\end{lemma}
\begin{proof}
Run monotone single-site uniform heat baths with common update variables
for pin values $-B,b,+B$. The two extreme chains can be maintained as
exact translates, since their conditional intervals are translates.
The middle chain remains between them: both endpoints of a conditional
interval are increasing functions of the neighboring heights. A stationary
subsequential coupling on the finite compact polytopes gives the asserted
coupling with the uniform marginals. Each tested statistic changes by at
most $B\sum|q_x|$. Condition first on mixed pin values and integrate the
same uniform bound.
\end{proof}

A string here is a prescribed set of positive or negative sites, each
having an oppositely prescribed neighbor. Thus its heights lie in $[-1,1]$.
Additional intervals $0\le h\le\eps$ on positive string sites are allowed.
After all string values are frozen, the remaining conditional law is
uniform on the corresponding Lipschitz polytope. Lemma~\ref{f:bounded}
therefore proves the moment assertion with $B=1$ on each component, provided
the anchor or zero-mass hypotheses are checked for the tested statistic.
A finite collection of components can instead be considered together
with the common pinned set. Cutting along pins and zero-extending the
zero-pinned comparison use the ambient lattice; neither operation requires
a regular shape for the boundary.

The same comparison transfers long-distance slack from zero pins to bounded
data at additive cost $2B$. At distances beyond a sufficiently large fixed
threshold this is absorbed into a relaxed linear slack. It does not
transfer microscopic strict gradients by itself.

\begin{lemma}[Small linear tilts]\label{f:tilts}
Under Proposition~\ref{f:averages}, let $Y=\sum_x b_xh(x)$, with support
in an $O(R)$ box and $|b_x|\le cR^{-2}$. Assume its box has the specified
connected pin anchor, or that $\sum_xb_x=0$. For sufficiently small fixed
$c>0$, the tilted density $e^Y/\E e^Y$ has uniformly bounded $L^2$ norm.
Hence for every event $A$,
\[
 \Prob_Y(A)\le C\Prob(A)^{1/2}.
\]
All fixed moments of admissible average statistics remain bounded under
this tilt, and any established arbitrary-parameter animal estimate is
inherited, with a harmless constant prefactor and changed parameter.
\end{lemma}
\begin{proof}
Apply \eqref{f:average-mgf} to $Y/c$ to bound $\E e^{2Y}$, and use
Jensen together with the uniform bound on $|\E Y|$ to bound $\E e^Y$
away from zero. Cauchy--Schwarz gives the assertions; for moments apply
it with the square of the desired moment. Multiple bounded tests can be
combined by Holder and a further reduction of $c$.
\end{proof}

\begin{lemma}[Strict filling with sign intervals]\label{f:strict-filling}
In a finite patch, prescribe arbitrary shell values, interior fixed pins
in $[-1/2,1/2]$, and intervals containing zero at any other prescribed
vertices. The intervals may be one-sided or of the form $[0,\eps]$.
Suppose a feasible filling $h$ satisfies
\begin{equation}\label{f:shell-room}
 |h(x)-b(s)|\le\operatorname{dist}(x,s)-2
\end{equation}
for every movable endpoint $x$ of a tested central edge and every shell
vertex $s$. If the edge is not fixed--fixed, there is a feasible filling
whose absolute gradient on that edge is at most $1/2$.
\end{lemma}
\begin{proof}
Recall the interval extension criterion. Data $[\alpha_z,\beta_z]$ at
specified vertices admit a 1-Lipschitz selection exactly when
$\alpha_z\le\beta_w+\operatorname{dist}(z,w)$ for every pair. Sufficiency
follows from the lower-cone function
$g(y)=\sup_z\{\alpha_z-\operatorname{dist}(z,y)\}$, which is Lipschitz,
lies above every lower endpoint and below every upper endpoint. Infinite
interval endpoints are omitted from the corresponding bound.

At a movable vertex $x$, the intervals propagated from all non-shell data
contain both $h(x)$ and zero. Indeed feasibility gives the first assertion.
For a nonfixed interval datum the second follows because its interval
contains zero. For a fixed pin $p$, its distance to $x$ is at least one,
whereas $|h(p)|\le1/2$, so its propagated interval contains zero as well.
The datum at $x$ itself, if present, also contains zero.

If both endpoints $u,v$ are movable, put
$c=\operatorname{median}(0,h(u),h(v))$ and prescribe $g(u)=g(v)=c$.
The value $c$ lies between zero and each old endpoint value; thus it
respects every propagated non-shell interval at both endpoints. Moreover
$|c-h(u)|,|c-h(v)|\le|h(u)-h(v)|\le1$.
Equation~\eqref{f:shell-room} preserves all endpoint--shell inequalities.
The new endpoint pair has difference zero, and every pair among the old
data was already consistent. The interval criterion therefore extends
the new prescriptions to a feasible filling.

If one endpoint is a fixed pin, prescribe zero at the other endpoint.
Its original height has absolute value at most $3/2$, so
\eqref{f:shell-room} again preserves all shell inequalities. The propagated
non-shell intervals contain zero, and the resulting edge gradient has
absolute value at most $1/2$. Apply the same extension criterion.
\end{proof}

\begin{corollary}[Sparse contacts with signs]\label{f:conditioned-sparse}
Fix a domain, fixed pin values in $[-1/2,1/2]$, and deterministic sign
strings as described above, allowing additional intervals $[0,\eps]$.
For every prescribed animal parameter there are a fixed box size and
$s>0$ such that, except on a bad set with that animal bound, the tested
central edges involving nonfixed vertices have $|\nabla h|\le1-s$.
String sites are included among the nonfixed vertices. Long-distance
slack and this edge slack persist with relaxed constants during a
sufficiently short stationary reflected step. Constants are uniform
over the strings' shape and over the positive interval lengths.
\end{corollary}
\begin{proof}
First freeze all string values, making them pins in $[-1,1]$.
Lemma~\ref{f:bounded} couples the conditional field within coordinatewise
distance one of the corresponding zero-pinned field. The controlled-gradient
witness argument gives arbitrary-parameter long-distance slack, with
the additive comparison loss absorbed by increasing its fixed threshold.
This bound is uniform over the conditioned string values, so integrating
preserves its joint form.

Now condition only on the exterior of each enlarged patch; its interior
strings remain variable. On a shell for which a filling has the preceding
long slack, central endpoints satisfy \eqref{f:shell-room}. Apply
Lemma~\ref{f:strict-filling}. If the conditional body has dimension $m$,
contraction toward its strict filling with ratio $1-2s$ sends the entire
body into $\{|\nabla h|\le1-s\}$. Therefore the conditional relative volume
of the near-contact region is at most
$1-(1-2s)^m\le2ms$. This remains true for arbitrarily thin sign intervals;
their effect on the body's total volume cancels in the relative-volume
estimate. Degenerate prescribed intervals are treated as fixed coordinates.

For disjoint enlarged patches the conditional laws are independent.
For a specified collection of bad boxes, either at least half have
long-slack failures, or at least half have near contacts with good long
slack. The joint bounds just proved apply to these two alternatives;
the subset count is absorbed by choosing the preliminary parameters
and then $s$ small enough. A finite coloring removes the separation
restriction. Lemma~\ref{f:step} gives the short-step statement. The tilt
transfer, when its statistic meets the anchored or zero-mass hypotheses,
is Lemma~\ref{f:tilts}.
\end{proof}

\begin{remark}[Fixed values after further reveals]\label{f:reveal-obligation}
The strict-filling proof uses the numerical
bound $1/2$ on fixed interior pins and uses intervals containing zero on
movable sign sites. After additional values are frozen, these hypotheses
must be checked again. In particular, arbitrary revealed values in
$[-1,1]$ cannot be substituted for the fixed $[-1/2,1/2]$ pins in
Lemma~\ref{f:strict-filling}. Uniform quenched contact sparsity under
such reveals is not asserted. The adapted comparisons below instead
use the annealed marginal-preservation argument of
Lemma~\ref{c:freezing}, with good marks recomputed after rejoining.
\end{remark}

\subsection{Distributional tightness and conditioning scope}
\begin{corollary}[Negative Sobolev tightness]\label{f:tightness}
Suppose the rescaled free sets lie in a fixed bounded box, their exterior
data are bounded, and a connected exterior pin set of macroscopic diameter
is available in a fixed larger box. The centered fields, made constant on
fundamental lattice cells of mesh $n^{-1}$ and extended by zero outside,
are tight in $H^{-s}$ on a fixed surrounding torus for every $s>0$.
\end{corollary}
\begin{proof}
Embed the lattice cells in a dyadic coordinate square with $2^J$ cells
on each side, where $2^J$ is comparable to $n$. Its physical size remains
bounded above and below. Expand the piecewise constant field in the
orthonormal two-dimensional Haar basis; the expansion terminates at the
cell scale. A Haar coefficient on a cell of physical side $r=R/n$ is
$r$ times a statistic satisfying \eqref{f:weights} and having exactly
zero total mass. For zero pins the tests include all ambient lattice
vertices, including the zero extension, so this cancellation also holds
across the boundary. Bounded data differ from the same zero-pinned sample
by a uniformly bounded field. Proposition~\ref{f:averages} consequently
gives a second moment at most $Cr^2$ for each centered coefficient.
The coarse scaling coefficient is bounded by the anchored assertion.

Write $F_j$ for the sum of wavelets at level $j$. There are $O(r_j^{-2})$
coefficients there, so $\E\|F_j\|_2^2\le C$. Each parent-cell integral
of $F_j$ vanishes. For $0<t\le1$, fractional Poincare on the parent cells
and duality give
\[
 \|F_j\|_{H^{-t}}\le Cr_j^t\|F_j\|_2.
\]
Minkowski in $L^2(\Prob;H^{-t})$ and the geometric sum of $r_j^t$ give
a uniform second moment in $H^{-t}$. Choose $0<t<s$ and use compactness
of its embedding into $H^{-s}$ on the fixed torus.
\end{proof}

For affine interpolation on full lattice triangles the same conclusion
holds. Its nodal basis and the fundamental-cell constant basis have equal
mass, bounded overlap, and supports at distance $O(n^{-1})$ from the node.
Testing their difference against $H^1$ functions gives an $H^{-1}$ error
at most $Cn^{-1}$ times the piecewise constant $L^2$ norm; interpolation
gives $Cn^{-t}$ for $0<t\le1$. Haar orthogonality in the preceding proof
gives $\E\|h-\E h\|_2^2\le C(1+\log n)$, so this error tends to zero
in every $H^{-t}$ with $t>0$. When interpolation is restricted to a polygon,
the boundary-cell contribution must also have vanishing area; this holds
for the stated smooth-domain approximation and bounded heights adjacent
to its fixed boundary.

Finally, all conditioning statements above specify the resulting kernel.
If a random exploration followed by freezing preserves a marginal law,
the corresponding annealed estimates are inherited by that marginal.
This does not give the same estimates uniformly for each realized revealed
set and every realized revealed value. Any later use of a quenched bound
must separately verify the fixed-value, geometric, and conditional-law
hypotheses of the result being invoked.
